\paragraph{Localizing names and transporting historical ink.}
The 96,000 training requests cover 6,000 shuffled plates from 1,985 families. Each plate crosses two characters with two stages and admits sixteen requests. The modern character, visible stage heading and attested raster occupy five 32-pixel cells per row. The model encodes a typed or pictured question into 24 image cells. A learned router distinguishes a plate from modern prose. With a plate, a recurrent reader scores the rows from the question, visible heading, character-image distance, step and previous emitted ink. The first fit trained 401,781 additional parameters for ten epochs. A subsequent stage-coverage comparison supervised overall presence.

The final matcher uses 132,960 parameters. For each contiguous question span $s$ of at most six cells and each visible header $h$, it computes $16(\langle f(s),g(h)\rangle-0.5)$ using normalized embeddings. Training compares genuine bilingual name spans with their headings and contrasts them with corrupted names and other training-question substrings. The highest scoring span supplies a location and a positive score marks a requested stage. The decoder draws the earliest outstanding stage, uses the frozen row scores within it, covers the emitted stage and stops when coverage is complete. The selected output contains all 1,024 physical bits from the chosen source cell. The eight training names and four headings constrain the interpretation of new prompt frames. At deployment the learned weights operate on visible cells and require no source database. On one RTX 4090 D the first fit took 38.19\,s and the final localization fit 3.21\,s, excluding the previous modern model, corpus preparation and evaluations. The transported pixels demonstrate source use within the trained geometry. Independent invention of an accurate historical form and generation of a factual explanation require separate evidence.
